\newpage

\section{Анализ динамики (2023--2025)}

Для анализа динамики использовался Датасет B, содержащий 4202 компании, присутствующие в обоих временных срезах.

\subsection{Матрица переходов статусов}

Матрица переходов (Рис.~\ref{fig:transition}) показывает, как изменились статусы компаний за два года. Основные наблюдения:

\begin{itemize}
    \item 2468 компаний сохранили статус active
    \item 233 компания перешла из active в acquired, 7 --- в public (итого 228 новых успешных случаев)
    \item 346 компаний перешли из active в inactive (закрылись)
    \item 627 компаний сохранили статус inactive, 499 --- acquired, 16 --- public
    \item Обратные переходы (из inactive в active или acquired) единичны (1 и 2 соответственно)
\end{itemize}

\begin{figure}[htbp]
    \centering
    \includegraphics[width=0.8\textwidth]{status_transition_matrix.png}
    \caption{Матрица переходов статусов (2023 -- 2025)}
    \label{fig:transition}
\end{figure}

Из матрицы переходов следует, что за два года из 3054 активных компаний 233 (7.6\%) достигли exit, а 346 (11.3\%) закрылись. Таким образом, соотношение <<закрытий>> к <<успехам>> составляет примерно 1.5:1, что согласуется с общей статистикой венчурного рынка.

\subsection{Динамика размера команды}

Анализ изменения размера команды (Рис.~\ref{fig:dynamics}) показал, что компании, достигшие успеха к 2025 году, чаще увеличивали размер команды по сравнению с неуспешными. Доля компаний, увеличивших команду, составляет около 11.1\% среди успешных и около 4.2\% среди неуспешных.

\begin{figure}[htbp]
    \centering
    \includegraphics[width=0.9\textwidth]{dynamics_team_size_delta.png}
    \caption{Динамика размера команды и успех (2023 -- 2025)}
    \label{fig:dynamics}
\end{figure}

Это наблюдение согласуется с результатами гипотезы H6 (найм как сигнал): рост команды является индикатором активного развития компании и коррелирует с достижением <<успеха>>. При этом важно отметить, что рост команды может быть как причиной, так и следствием успеха (привлечение инвестиций позволяет нанимать больше сотрудников).
